\documentclass[10pt]{beamer}

\usepackage[latin1]{inputenc}
\usepackage[spanish]{babel}

\mode<presentation>
\usetheme{Madrid}
\usecolortheme[RGB={128,0,0}]{structure}
\usefonttheme{serif}
\useinnertheme{rectangles}
\useoutertheme{split}
\setbeamercovered{transparent}

\title{Python funcional}
\author{Jes�s Espino Garc�a}
\date{8 de Noviembre de 2014}
\subject{Python funcional}

\institute[Kaleidos \& Taiga.io]{
    \includegraphics[height=1.5cm]{kaleidos}
    \includegraphics[height=1.5cm]{taigaio} \\
    \includegraphics[height=1.5cm]{pyladies}
}

\begin{document}

  \frame{\maketitle}
  
  \section{Introducci�n}

  \begin{frame}{Introducci�n}
    \begin{center}Introducci�n\end{center}
  \end{frame}
  
  \begin{frame}{�Que significa funcional?}
    \begin{itemize}
      \item Programaci�n con funciones (matematicas)
      \item Funciones puras (mismas entradas, mismas salidas).
      \item Inmutabilidad.
      \item Ausencia de estado.
    \end{itemize}
  \end{frame}

  \begin{frame}{�Por qu�?}
    \begin{itemize}
      \item Paralelizaci�n.
      \item Facilitar el testing.
      \item Reusabilidad.
      \item Depuraci�n m�s f�cil.
      \item Estado muy controlado.
    \end{itemize}
  \end{frame}
  
  \begin{frame}{Tipicas estrategias funcionales}
    \begin{itemize}
      \item Combinaci�n y composici�n de funciones peque�as.
      \item Datos + funciones transformadoras.
      \item Aplicaci�n de transformaciones mediante orden superior.
      \item Uso de funciones inline.
      \item Acotado de efectos laterales.
      \item Tendencia al uso de funciones puras.
    \end{itemize}
  \end{frame}

  \begin{frame}{�Qu� necesito?}
    \begin{itemize}
      \item Funciones como ciudadanos de primera (son un objeto m�s).
    \end{itemize}
  \end{frame}

  \begin{frame}{�Es python un lenguaje funcional?}
    \begin{itemize}
      \item No.
      \item Es un lenguaje multi-paradigma.
      \item Soporta algunas caracteristicas funcionales.
      \item Permite hacer programaci�n funcional.
      \item Carece de caracteristicas avanzadas presentes en lenguajes funcionales.
    \end{itemize}
  \end{frame}

  \begin{frame}{�Que me dan los lenguajes funcionales?}
    \begin{itemize}
      \item Estructuras inmutables eficientes.
      \item Funciones de orden superior.
      \item Pattern matching.
      \item TCO: Tail call optimization.
      \item Aplicaci�n parcial y currificaci�n.
      \item Control de efectos laterales.
      \item Funciones lambda.
      \item Evaluaci�n peresoza.
      \item Composici�n de funciones.
    \end{itemize}
  \end{frame}

  \begin{frame}{�Que me da python?}
    \begin{itemize}
      \item Evaluaci�n peresoza (Limitada).
      \item Aplicaci�n parcial.
      \item Funciones lambda.
      \item Funciones de orden superior.
    \end{itemize}
  \end{frame}

  \begin{frame}{�Que me da fn.py?}
    \begin{itemize}
      \item Estructuras inmutables eficientes (En desarrollo).
      \item TCO: Tail call optimization.
      \item Aplicaci�n parcial y currificaci�n.
      \item Composici�n de funciones.
      \item Funciones lambda (Al estilo de Scala).
    \end{itemize}
  \end{frame}

  \begin{frame}[containsverbatim]{Funcional vs. Imperativo}
    \begin{block}{Imperativo}
       \begin{verbatim}
x = sum(1, 2)
y = sum(x, 3)
z = prod(y, 4)
print(z)
       \end{verbatim}
    \end{block}
    \begin{block}{Funcional}
       \begin{verbatim}
print(prod(sum(sum(1,2),3), 4))
       \end{verbatim}
    \end{block}
  \end{frame}

  \begin{frame}[containsverbatim]{Funcional vs. Imperativo}
    \begin{block}{Funcional con composici�n y aplicacion parcial}
       \begin{verbatim}
func = F(sum, 1, 2) >> F(sum, 3) >> F(prod, 4) >> print
func()
       \end{verbatim}
    \end{block}
    \begin{block}{Funcional con currificaci�n}
       \begin{verbatim}
prod4 = prod(4)
sum3 = sum(3)
sum2 = sum(2)
print(prod4(sum3(sum2(1))))
       \end{verbatim}
    \end{block}
  \end{frame}

  \section{Python funcional}

  \begin{frame}{Python funcional}
    \begin{center}Python funcional\end{center}
  \end{frame}

  \begin{frame}{Evaluaci�n peresoza}
    \begin{itemize}
      \item Iteradores
      \item Generadores
    \end{itemize}
  \end{frame}

  \begin{frame}[containsverbatim]{Evaluaci�n peresoza}
    \begin{block}{Iteradores}
       \begin{verbatim}
>>> i = map(print, [1,2,3])
>>> next(i)
1
>>> i = map(print, [1,2,3])
>>> list(i)
1
2
3
[None, None, None]
       \end{verbatim}
    \end{block}
  \end{frame}

  \begin{frame}[containsverbatim]{Evaluaci�n peresoza}
    \begin{block}{Generadores}
       \begin{verbatim}
>>> import itertools
>>> def generate():
...    x = 0
...    while True:
...        yield x
...        x += 1
>>> numbers = generate()
>>> list(itertools.takewhile(lambda x: x < 10, numbers))
[0, 1, 2, 3, 4, 5, 6, 7, 8, 9]
>>> list(itertools.takewhile(lambda x: x < 12, numbers))
[11]
       \end{verbatim}
    \end{block}
  \end{frame}

  \begin{frame}[containsverbatim]{Aplicaci�n parcial}
    \begin{block}{Applicaci�n parcial}
       \begin{verbatim}
>>> from functools import partial
>>> from opertor import add
>>> add4 = partial(add, 4)
>>> add4(3)
7
>>> print_noline = partial(print, end="")
>>> print_noline("hello")
hello>>>
       \end{verbatim}
    \end{block}
  \end{frame}

  \begin{frame}[containsverbatim]{Funciones lambda}
    \begin{block}{Funciones lambda}
       \begin{verbatim}
>>> pow2 = lambda x: x**2
>>> pow2(10)
100
       \end{verbatim}
    \end{block}
  \end{frame}

  \begin{frame}{Funciones de orden superior}
    \begin{itemize}
      \item map
      \item filter
      \item sorted
      \item reduce
      \item decorators
    \end{itemize}
  \end{frame}

  \begin{frame}[containsverbatim]{Funciones de orden superior}
    \begin{block}{Funciones de orden superior}
       \begin{verbatim}
>>> list(map(lambda x: x**2, [1, 2, 3]))
[1, 4, 9]
>>> list(filter(lambda x: x > 1, [1, 2, 3]))
[2, 3]
>>> sorted([2, 1, 3], key=lambda x: x)
[1, 2, 3]
>>> sorted([1, 2, 3], key=lambda x: -x)
[3, 2, 1]
>>> from functools import reduce
>>> reduce(lambda x, y: x + y, [1, 2, 3])
6
>>> from functools import lru_cache
>>> cached_sum = lru_cache()(lambda x: sum(range(x)))
>>> cached_sum(4)
6
       \end{verbatim}
    \end{block}
  \end{frame}

  \section{fn.py}

  \begin{frame}{fn.py}
    \begin{center}fn.py\end{center}
  \end{frame}

  \begin{frame}{Estructuras inmutables}
    \begin{itemize}
      \item LinkedList
      \item Stack
      \item Queue
      \item Deque
      \item Vector
      \item SkewHeap
      \item PairingHeap
    \end{itemize}
  \end{frame}

  \begin{frame}[containsverbatim]{Estructuras inmutables}
    \begin{block}{LinkedList}
      \begin{verbatim}
>>> from fn.immutable import LinkedList
>>> l = LinkedList()
>>> l.head
>>> l.tail
>>> l2 = l.cons(10)
>>> l2.head
10
>>> l2.tail
<fn.immutable.list.LinkedList object at 0x7f3927e59f08>
>>> l.head
>>> l.tail
      \end{verbatim}
    \end{block}
  \end{frame}

  \begin{frame}[containsverbatim]{Estructuras inmutables}
    \begin{block}{Stack}
      \begin{verbatim}
>>> from fn.immutable import Stack
>>> s = Stack()
>>> s.head
>>> s.tail
>>> s2 = s.push(10)
>>> s2.head
10
>>> s2.tail
<fn.immutable.list.Stack object at 0x7f3926ae9818>
>>> (value, s3) = s2.pop()
>>> value
10
      \end{verbatim}
    \end{block}
  \end{frame}


  \begin{frame}[containsverbatim]{TCO}
    \begin{block}{Normal recursion}
      \begin{verbatim}
def fact(n):
    if n == 0: return 1
        return n * fact(n-1)
      \end{verbatim}
    \end{block}

    \begin{block}{TCO}
      \begin{verbatim}
from fn import recur

@recur.tco
def fact(n, acc=1):
    if n == 0: return False, acc
    return True, (n-1, acc*n)
      \end{verbatim}
    \end{block}
  \end{frame}

  \begin{frame}[containsverbatim]{Applicaci�n parcial}
    \begin{block}{Applicaci�n parcial}
      \begin{verbatim}
>>> from fn import F
>>> from operator import add
>>> add2 = F(add, 2)
>>> add2(3)
5
      \end{verbatim}
    \end{block}
  \end{frame}

  \begin{frame}[containsverbatim]{Currificacion}
    \begin{block}{Currificacion}
      \begin{verbatim}
>>> from fn.func import curried
>>> curry_add = curried(lambda x, y: x + y)
>>> curry_add(2)(3)
5
>>> @curried
... def curried_add(x, y):
...     return x + y
...
>>> curried_add(2)(3)
5
      \end{verbatim}
    \end{block}
  \end{frame}
  
  \begin{frame}[containsverbatim]{Composici�n de funciones}
    \begin{block}{Composici�n normal}
      \begin{verbatim}
>>> myfunc = lambda x: duplicate(add2(x))
>>> myfunc(3)
10
      \end{verbatim}
    \end{block}

    \begin{block}{Composici�n al estilo fn.py}
      \begin{verbatim}
>>> myfunc = F(duplicate) << add2
>>> myfunc(3)
10
>>> myfunc = F(add2) >> duplicate
>>> myfunc(3)
10
      \end{verbatim}
    \end{block}
  \end{frame}

  \begin{frame}[containsverbatim]{Funciones lambda al estilo scala}
    \begin{block}{Funciones lambda al estilo scala}
      \begin{verbatim}
>>> from fn import _
>>> (_ + _)(2, 3)
5
>>> list(map(_ + 2, [1, 2, 3]))
[3, 4, 5]
      \end{verbatim}
    \end{block}
  \end{frame}

  \section{Para terminar}

  \begin{frame}{Para terminar}
    \begin{center}Para terminar\end{center}
  \end{frame}

  \begin{frame}{Consclusiones}
    \begin{itemize}
      \item Python permite programar de forma funcional.
      \item Fn.py nos da las herramientas para llegar un poco m�s lejos.
      \item Python + Fn.py se queda lejos de lenguajes como Erlang, Clojure o Haskell.
      \item Lo que se puede aplicar en Python es una mejora significativa sobre el codigo.
    \end{itemize}
  \end{frame}

  \begin{frame}{Referencias}
    \begin{itemize}
      \item https://github.com/kachayev/fn.py: Fn.py
      \item https://docs.python.org/3/howto/functional.html: Howto de programaci�n funcional.
      \item http://kachayev.github.io/talks/uapycon2012/: Charla de Alexey Kachayev
    \end{itemize}
  \end{frame}

  \begin{frame}{Dudas}
    \dots
  \end{frame}

\end{document}
